% !TEX root = ../aa_SoM_Chapters.tex

% ö = \%c3\%b6
% ä = \%C3\%A4

\xdef\sectiontitle{\Isectiontitle} %aiheuttama

%%%%%%%%%%%%%%%
\begin{frame}[label=1_2_a]%[label=1_2_TOC1]
\frametitle{\texorpdfstring{\culine[\sectiontitleulinecolor]{\sectiontitle}}{\sectiontitle} }

\vspace{0.5cm}

\begin{itemize}[leftmargin=0.5cm,itemsep=2mm]
    \item[1.1]<1-> \hyperlink{1_1_a}{\IaSubsectiontitle}
    \item[1.2]<1-> \hyperlink{1_2_a}{\Alert[red]{\IbSubsectiontitle}}	
    \item[1.3]<1-> \hyperlink{1_3_a}{\IcSubsectiontitle}	
    \item[1.4]<1-> \hyperlink{1_4_a}{\IdSubsectiontitle}
    \item[1.5]<1-> \hyperlink{1_5_a}{\IeSubsectiontitle}
    \item[1.6]<1-> \hyperlink{1_6_a}{\IfSubsectiontitle}
    \item[1.7]<1-> \hyperlink{1_7_a}{\IgSubsectiontitle}
\end{itemize}

\vspace{-1cm}
\begin{itemize}[leftmargin=9.5cm,itemsep=2mm]
    \item[1.8]<1-> \hyperlink{1_8_a}{\IhSubsectiontitle}
	\item[1.9]<1-> \hyperlink{1_9_a}{\IiSubsectiontitle}
	\item[1.10]<1-> \hyperlink{1_10_a}{\IjSubsectiontitle}
	\item[1.11]<1-> \hyperlink{1_11_a}{\IkSubsectiontitle}
\end{itemize}

%\endofslide 
\end{frame}
%%%%%%%%%%%%%%%

%\xdef\IbSubsectiontitle{Entropy and Temperature}
\xdef\sectiontitle{\IbSubsectiontitle} 

%%%%%%%%%%%%%%%
\begin{frame}
\frametitle{\texorpdfstring{\culine[\sectiontitleulinecolor]{\sectiontitle}}{\sectiontitle} }

\begin{itemize}[itemsep=1.9mm]
	\item \CDAlert[red]{Entropy $S$} is a quantity proportional to the logarithm of the \CDAlert[tunired]{'number of ways' $W$} \appear{(i.\,e. to the log of \CDAlert[red]{multiplicity $\Omega$}).} \appear{Boltzmann constant $k_B$ sets the dimensions for the entropy $S$}
	\[
S  \equal \appear{k_{B} \ln W} \qquad \appear{\&} \qquad \appear{S}  \equal  \appear{k_{B} \ln \Omega} \qquad \appear{\&} \qquad \appear{[S]=\frac{J}{K}}
\]

\item We saw that the system has a tendency to \CDAlert[blue]{evolve towards} \CDAlert[blue]{equilibrium}\appear{, i.\,e. towards highest $W$}

\item According to the \CDAlert[red]{second law of thermodynamics}\appear{,  entropy tends to increase (in a closed system): }
\[
\Delta S \appear{\geq 0}
\]
\appear{Considering the number of ways $W$:} 
\[
\Delta S  \equal \appear{S_{f}}  \minus \appear{S_{i}}  \equal \appear{k_{B} \textrm{[[} } \appear{\ln W_{f}} \appear{-\ln W_{i} \textrm{]]}}  \equal \appear{k_{B} \ln \frac{W_{f}}{W_{i}} }
\]

\end{itemize}

\endofslide 
%\endofsection
\end{frame}
%%%%%%%%%%%%%%%

%%%%%%%%%%%%%%%
\begin{frame}
\frametitle{\texorpdfstring{\culine[\sectiontitleulinecolor]{\sectiontitle}}{\sectiontitle} }

\appear{\footnote{{Recap the definitions for isolated, \Alert[red]{closed}, and open systems.}}} 
\begin{minipage}{8.5cm}
\begin{itemize}
	\item Consider two bodies\appear{, 1 and 2}\appear{, in \CDAlert[red]{thermal} \CDAlert[red]{contact}.} \appear{Heat will flow, and entropy will change:}
\[ 
 d S_{t o t} \equal \appear{d S_{1} } \appear{+ d S_{2}} 
 \]
\appear{Assume that entropy changes due to changes in the energy} \appear{(ok, since here, only heat}
\end{itemize}
\end{minipage}
\vspace{-0mm}
\begin{itemize}
\item[] flows), so $S=S(E)$\appear{, and therefore}
\[ 
 d S \equal  \appear{\bigg(} \frac{\partial S}{\partial E} \appear{\bigg)} \,\appear{d E} 
 \]
\item Thus $\appear{d S_{t o t} =} \appear{\textrm{((} \frac{\partial S}{\partial E}} \appear{\textrm{))}_{2}} \appear{d E_{2}}  \plus  \appear{\textrm{((} } \frac{\partial S}{\partial E} \appear{\textrm{))}_{1}} \appear{d E_{1}}$\appear{,
and since energy is transferred only between 1 and 2}\appear{, condition $d E_{1}=-d E_{2}$ applies leading to}
\[ 
d S_{tot} \equal   \LB   \appear{\bigg( \frac{\partial S}{\partial E}  \bigg)_{2}} \minus  \appear{\bigg( \frac{\partial S}{\partial E}  \bigg)_{1}}  \;\RB \appear{d E_{2}}
 \]
\end{itemize}

\xdef\loc{north east}
\begin{tikzpicture}[remember picture,overlay] % anchor=south
    \node (A) [yshift=1*\figYshift,xshift=\figXshift, anchor=\loc] at (current page.\loc) {\includegraphics[page=1,width=4.8cm]{Figures_booklet/SOM_9_Statistical_mechanics_figures/SOM_Figure_4.png}}; %scale=\figscale. % 8cm max hei
\end{tikzpicture}

\endofslide 
%\endofsection
\end{frame}
%%%%%%%%%%%%%%%

%%%%%%%%%%%%%%%
\begin{frame}
\frametitle{\texorpdfstring{\culine[\sectiontitleulinecolor]{\sectiontitle}}{\sectiontitle} }

\begin{itemize}
	\item \CDAlert[tunired]{At equilibrium the two bodies will have the same temperature}

\item We also know that when equilibrium is reached\appear{,  the entropy $S$ does not increase anymore}\appear{, so $d S_{t o t}=0$ resulting in} 
\[ 
  \bigg(  \Frac{\partial S}{\partial E}  \bigg)_{2} \equal  \appear{ \bigg( \Frac{\partial S}{\partial E}  \bigg)_{1}} 
 \]
\item The above condition leads to a convenient definition for the temperature: 

\item[] \begin{itemize}
	\item At equilibrium the temperatures are equal $\appear{T_1}  \equal  \appear{ T_2} \equiv \appear{T}$
	\item We can \CDAlert[red]{define temperature $T$} via equation 
	 \[
	 \frac{1}{T} \equiv \appear{\bigg(} \frac{\partial S}{\partial E} \appear{\bigg)\,,} \qquad \appear{\,\textrm{~where}} \qquad \appear{S}  \equal \appear{k_{B} \ln W}
	 \]
\end{itemize}
\end{itemize}
 
 
\endofslide 
%\endofsection
\end{frame}
%%%%%%%%%%%%%%%

%%%%%%%%%%%%%%%
\begin{frame}
\frametitle{\texorpdfstring{\culine[\sectiontitleulinecolor]{\sectiontitle}}{\sectiontitle} }

\begin{itemize}
	
\item \CDAlert[blue]{How can we take a derivative of $S$ with respect to $E$?} 
\begin{itemize}
	\item In a physical system usually also the number of possibilities \appear{(number of 'ways' $W$)} \appear{is a function of energy} \appear{(more energy, higher multiplicity)} \appear{so there is no contradiction}
\end{itemize}
 

\item In many cases $W=W(E)$ \appear{and thus also the entropy $S=S(E)$}
\item Recalling that the change in energy $dE$ \appear{is equal to the heat transferred between objects} \appear{($dE=dQ$)}\appear{,  we get }
\[
	 \Frac{1}{T}  \equal   \frac{\partial S}{\partial E}  \qquad \Rightarrow \qquad \appear{dS}  \equal  \frac{dQ}{T}  
	 \]
\appear{allowing us to calculate the \CDAlert[red]{change of entropy $\Delta S$} of the system:}
\[ 
 \Delta S  \equal \appear{S_{2}}  \minus  \appear{S_{1}} \equal \appear{\nint_{i}^{f}} \frac{d Q}{T} 
 \]

\end{itemize}
 

%\endofslide 
\endofsection
\end{frame}
%%%%%%%%%%%%%%%

%%%%%%%%%%%%%%%%
%\begin{frame}
%\frametitle{\texorpdfstring{\culine[\sectiontitleulinecolor]{\sectiontitle}}{\sectiontitle} }
%
%\begin{itemize}
%	\item 
%\end{itemize}
%
%%\xdef\loc{north east}
%%\begin{tikzpicture}[remember picture,overlay] % anchor=south
%%    \node (A) [yshift=1*\figYshift,xshift=\figXshift, anchor=\loc] at (current page.\loc) {\includegraphics[page=1,width=4.8cm]{Figures_booklet/SOM_9_Statistical_mechanics_figures/SOM_Figure_4.png}}; %scale=\figscale. % 8cm max hei
%%\end{tikzpicture}
%
%\endofslide 
%%\endofsection
%\end{frame}
%%%%%%%%%%%%%%%%

